\section{固定模型实验：反馈、有效候选与搜索效率}\label{sec:feedback}
\subsection{最终结果与候选质量}
固定模型实验把生成模型统一为托管 Qwen3.8-27B，用三种信息条件观察搜索过程。ahc024 的反馈与独立探索各运行三次；ahc026 进一步加入无后续反馈组。实验共提出 96 个候选，60 个通过开发验证。最终交付 15 个程序，均通过对应的 15 个最终用例。

\reportfigure[0.92\linewidth]{feedback_ablation.pdf}{各组的最终得分相对初始程序的提升。圆点展示三次重复，横线表示中位数；空心点为交付初始程序，叉号为接口错误终止后保留正确候选的运行。两道题使用各自纵轴刻度。}{fig:feedback}

\begin{smalltable}
\begin{tabularx}{\linewidth}{@{}llrrrL@{}}\toprule
任务 & 策略 & 最终中位分 & 提升 & 候选通过 & 初始程序交付\\\midrule
ahc024 & F & 15 & 0.000\% & 5/25 & 3/3\\
ahc024 & I & 15 & 0.000\% & 15/23 & 3/3\\
ahc026 & F & 132273 & 0.310\% & 14/20 & 1/3\\
ahc026 & I & 133031 & 0.885\% & 10/12 & 1/3\\
ahc026 & N & 132428 & 0.428\% & 16/16 & 0/3\\\bottomrule
\end{tabularx}
\end{smalltable}

两道题体现了不同的难点。ahc024 的六轮都停留在初始得分，模型首先需要解决地图连通性与邻接关系的联合约束。ahc026 三组都产生了高于初始程序的候选，最终中位数依次为 I、N、F；其中 I/r2 的研究因 HTTP 400 结束，已有候选仍取得 135088 分。该候选与失败状态共同保留在统计中。

\subsection{从通过验证到刷新最佳}
一次正确候选可以优于初始程序，也可以与初始程序等分；即使优于初始程序，它仍可能低于本轮已经找到的最佳方案。为观察搜索效率，图 \ref{fig:outcomes} 先将候选划为三个互斥类别：正确且优于初始、正确但未优于初始、验证失败。每轮保留最佳的变化另由过程轨迹描述。

\reportfigure[0.80\linewidth]{candidate_outcomes.pdf}{各组生成候选的结果构成。是否优于初始程序使用开发集裁判分数判断；“优于初始”与“刷新当前最佳”对应不同的搜索阶段。}{fig:outcomes}

ahc024 反馈组的 25 个候选中，15 个输出违反任务约束，5 个发生运行错误，5 个正确但得分没有增加。独立探索组的 23 个候选中，7 个输出错误、1 个运行错误，其余 15 个保持正确但无增益。两组分别消耗 56 与 26 次模型请求；反馈组包含 28 次分析动作，独立组的 23 次提案直接从初始程序出发。

在 ahc026 上，反馈组的 20 个候选包含 9 个优于初始、5 个正确无初始增益、6 个运行错误；独立组对应为 5、5、2；无后续反馈组为 9、7、0。反馈组 20 个候选只有 9 种精确代码哈希，说明部分预算再次落在已出现的实现上。这个现象把改进方向指向候选去重、失败原因利用和实验动作的具体约束。

\subsection{失败案例：地图压缩中的连通性约束}
ahc024 反馈组第一轮的研究循环用时 328.337 秒，共 19 次请求：10 次分析、8 次提案和 1 次结束。八个候选中，只有 c4 通过验证。

\begin{smalltable}\renewcommand{\arraystretch}{1.12}%
\begin{tabularx}{\linewidth}{@{}lLL@{}}\toprule
候选 & 尝试及运行事实 & 对后续研究的含义\\\midrule
c1 & 程序引用未定义变量 \texttt{color}，触发 NameError & 基础执行错误阻断了方案验证\\
c2、c3 & 通过保留边界、连接区域等方式减少格子；裁判判为非法输出 & 需要同时维护区域连通性与原有邻接关系\\
c4 & 从边界格子向同色区域执行遍历，最终覆盖原区域全部格子；五例均正确，合计 5 分 & 恢复合法性，同时也恢复了原始地图规模\\
c5、c6 & 两份候选代码哈希相同，均为非法输出 & 重复提案消耗请求与评测时间\\
c7、c8 & 后续删格策略仍违反任务约束 & 本轮交付选择器保留 c0\\\bottomrule
\end{tabularx}
\end{smalltable}

从地图内部删去一个格子，会把它变为颜色 0。若新空白与外部空白无法连通，输出就会失去合法性；删除同色区域的连接点，还可能切断该颜色的连通分量。有效操作需要把这些约束变为可检查条件。该轮围绕整图构造与修补进行调整，始终没有形成既合法又能增加空白格的操作序列。

ahc026 的反馈组最终中位分低于独立探索组；ahc024 两组均交付初始程序。提升反馈利用效率，需要让错误信息落实到可执行检查、局部变换和下一次实验目标上。
\source{[S2] 原始候选与逐轮报告；候选分类和重复哈希由分析脚本重算。}
